\section{31行业预告与阶段全景}

行业表先回答“哪里有盈利变化”，再回答“价格和资金是否已经确认”。有色金属、电子、非银金融、基础化工和电力设备是预告密度最高的五个行业，但高影响数量多并不等于低位布局：电子和基础化工在本轮行业资金模型中仍是无信号，有色金属多数只是单日反弹。

\small
\begin{longtable}{L{1.8cm}L{1.55cm}R{1.05cm}R{1.05cm}R{1.05cm}R{1.25cm}L{4.1cm}}
\toprule
\textbf{行业} & \textbf{阶段} & \textbf{预告数} & \textbf{正利润} & \textbf{高影响} & \textbf{H1净利和} & \textbf{代表公司} \\
\midrule
\endfirsthead
\toprule
\textbf{行业} & \textbf{阶段} & \textbf{预告数} & \textbf{正利润} & \textbf{高影响} & \textbf{H1净利和} & \textbf{代表公司} \\
\midrule
\endhead
有色金属 & 单日反弹 & 25 & 25 & 16 & 1071.9 & 洛阳钼业、宏桥控股、紫金矿业 \\
电子 & 无信号 & 45 & 41 & 10 & 1231.3 & 长鑫科技、江波龙、兆易创新 \\
非银金融 & 低位无资金 & 14 & 14 & 10 & 777.4 & 招商证券、中信证券、中金公司 \\
基础化工 & 无信号 & 51 & 46 & 7 & 395.1 & 万华化学、卫星化学、盐湖股份 \\
电力设备 & 无信号 & 28 & 24 & 6 & 173.6 & 天赐材料、天华新能、恩捷股份 \\
交通运输 & 单日反弹 & 8 & 8 & 4 & 170.7 & 招商轮船、海通发展、中远海能 \\
计算机 & 资金观察 & 15 & 12 & 4 & 61.3 & 浪潮信息、智微智能、紫光股份 \\
石油石化 & 资金观察 & 12 & 11 & 3 & 195.8 & 恒逸石化、东方盛虹、恒力石化 \\
机械设备 & 单日反弹 & 27 & 20 & 3 & 108.0 & 大族数控、东方精工、大族激光 \\
国防军工 & 启动确认 & 6 & 4 & 2 & 59.0 & 松发股份、高德红外 \\
医药生物 & 单日反弹 & 18 & 14 & 2 & 40.2 & 海思科、海正药业 \\
煤炭 & 静默吸纳 & 4 & 2 & 1 & 115.2 & 陕西煤业 \\
公用事业 & 单日反弹 & 8 & 8 & 1 & 53.1 & 华润新能源 \\
汽车 & 单日反弹 & 16 & 12 & 1 & 28.2 & 宁波华翔 \\
环保 & 单日反弹 & 5 & 5 & 1 & 25.8 & 浙富控股 \\
传媒 & 启动确认 & 6 & 3 & 1 & 23.8 & 巨人网络 \\
农林牧渔 & 启动确认 & 7 & 2 & 1 & 3.9 & 益生股份 \\
食品饮料 & 单日反弹 & 11 & 10 & 0 & 36.6 & --- \\
通信 & 无信号 & 8 & 8 & 0 & 19.4 & --- \\
建筑材料 & 无信号 & 4 & 2 & 0 & 17.1 & --- \\
纺织服饰 & 单日反弹 & 7 & 6 & 0 & 11.8 & --- \\
建筑装饰 & 低位无资金 & 6 & 4 & 0 & 8.7 & --- \\
社会服务 & 低位无资金 & 5 & 2 & 0 & 8.7 & --- \\
轻工制造 & 单日反弹 & 9 & 6 & 0 & 8.5 & --- \\
商贸零售 & 单日反弹 & 5 & 4 & 0 & 6.2 & --- \\
钢铁 & 单日反弹 & 5 & 3 & 0 & 4.5 & --- \\
家用电器 & 单日反弹 & 3 & 1 & 0 & 1.6 & --- \\
房地产 & 静默吸纳 & 6 & 1 & 0 & 0.9 & --- \\
银行 & 静默吸纳 & 0 & 0 & 0 & 0.0 & --- \\
综合 & 无信号 & 0 & 0 & 0 & 0.0 & --- \\
美容护理 & 低位无资金 & 0 & 0 & 0 & 0.0 & --- \\
\bottomrule
\end{longtable}
\normalsize
\sourcenote{data/full\_market\_preview\_screen\_20260712.json；利润单位为亿元。}

\section{板块判断：静默、低位与已启动}

\begin{exhibitbox}[Exhibit 6：板块机会地图]
\begin{tabularx}{\textwidth}{L{2.2cm}L{3.0cm}L{3.4cm}X}
\toprule
\textbf{类型} & \textbf{行业} & \textbf{核心证据} & \textbf{投资含义} \\
\midrule
静默吸纳 & 煤炭、银行、房地产 & 行业阶段低，煤炭出现陕西煤业高影响预告 & 行业方向成立，但正式目标仍需逐票通过现价估值 \\
低位盈利 & 有色、石化、非银、汽车、电力设备、计算机 & 11只优先股位于一年区间约9\%--33\% & 赔率来自位置与盈利改善，主要风险是缺少持续资金确认 \\
已启动仍有空间 & 石化、军工、传媒、创新药、AI基础设施 & 趋势确认与盈利桥并存 & 只做回撤和业绩确认，不能把同比高增线性外推 \\
业绩兑现但价格先行 & 电子、计算机、军工等 & 江波龙、兆易创新、香农芯创等涨幅或位置较高 & 等待新分母、现金流和回撤，不追逐旧目标价 \\
\bottomrule
\end{tabularx}
\end{exhibitbox}

差异化结论是：本轮最广泛的盈利改善来自有色、电子和非银金融，但真正符合“尚未充分启动”的是个股而非整板块；石油石化同时拥有低位的恒力石化和已启动的恒逸石化、东方盛虹，最适合用左右两侧篮子拆分，而不是统一看多。
